\section{Propositional Logic}
\label{sec:propositional-logic}

\subsection{Basic Logical Equivalences}


\subsubsection{Commutative Law}
\begin{itemize}
  \item $p \wedge q \equiv q \wedge p$
  \item $p \vee q \equiv q \vee p$
\end{itemize}

\subsubsection{Associative Law}
\begin{itemize}
  \item $(p \wedge q) \wedge r \equiv p \wedge (q \wedge r)$
  \item $(p \vee q) \vee r \equiv p \vee (q \vee r)$
\end{itemize}

\subsubsection{Distributive Law}
\begin{itemize}
  \item $p \wedge (q \vee r) \equiv (p \wedge q) \vee (p \wedge r)$
  \item $p \vee (q \wedge r) \equiv (p \vee q) \wedge (p \vee r)$
\end{itemize}

\subsubsection{Identity Law}
\begin{itemize}
  \item $p \wedge \taut \equiv p$
  \item $p \vee \cont \equiv p$
\end{itemize}

\subsubsection{Negation Law}
\begin{itemize}
  \item $p \vee \negt p \equiv \taut$
  \item $p \wedge \negt p \equiv \cont$
\end{itemize}

\subsubsection{Double Negation Law}
\begin{itemize}
  \item $\negt (\negt p) \equiv p$
\end{itemize}

\subsubsection{Idempotent Law}
\begin{itemize}
  \item $p \wedge p \equiv p$
  \item $p \vee p \equiv p$
\end{itemize}

\subsubsection{Universal Bound Law}
\begin{itemize}
  \item $p \vee \taut \equiv \taut$
  \item $p \wedge \cont \equiv \cont$
\end{itemize}

\subsubsection{De Morgan's Law}
\begin{itemize}
  \item $\negt (p \wedge q) \equiv \negt p \vee \negt q$
  \item $\negt (p \vee q) \equiv \negt \wedge \negt q$
\end{itemize}

\subsubsection{Absorption Law}
\begin{itemize}
  \item $p \vee (p \wedge q) \equiv p$
  \item $p \wedge (p \vee q) \equiv p$
\end{itemize}

\subsubsection{Negations of \texorpdfstring{$\taut$}{t} and \texorpdfstring{$\cont$}{c}}
\begin{itemize}
  \item $\negt \taut \equiv \cont$
  \item $\negt \cont \equiv \taut$
\end{itemize}


\subsection{XOR}
\begin{itemize}
  \item $p \oplus q \equiv (p \vee q) \wedge \negt (p \wedge q)$
  \item $p \oplus q \equiv (a \wedge \negt b) \vee (\negt a \wedge b)$
\end{itemize}


\subsection{Implication (If-Then)}
\begin{itemize}
  \item $p \rightarrow q \equiv \negt p \vee q$
\end{itemize}

\subsection{Equivalence (Biconditional If and Only if)}
\begin{itemize}
  \item $p \leftrightarrow q \equiv (p \rightarrow q) \wedge (q \rightarrow p)$
\end{itemize}


\section{Digital Logic Circuits}
\label{sec:digital-logic-circuits}

\subsection{Disjunctive Normal Form (DNF)}

"OR of ANDs"

\textbf{From truth table to DNF}
\begin{enumerate}
  \item Identify the rows whose output is $1$
  \item AND together all variables in the same row (negate the variable when $0$)
  \item OR together the $1$-rows
\end{enumerate}

\subsection{Conjunctive Normal Form (CNF)}

"AND of ORs"

\textbf{From truth table to CNF}
\begin{enumerate}
  \item Identify the rows whose output is $0$
  \item OR together all variables in the same row (negate the variable when $1$)
  \item AND together the $0$-rows
\end{enumerate}

\subsection{NOR Gate}

\subsubsection{NOT Using NOR}

\begin{circuitikz}
  \draw
    node[nor port] (nor) {}

    % Label the output
    (nor.out) node[right] {$\negt a$}

    % Tie both inputs together with a wire and label
    (nor.in 1) coordinate (top)
    (nor.in 2) coordinate (bot)
    (top) -- (bot)

    % Single input wire from the midpoint
    ($(top)!0.5!(bot)$) -- ++(-0.5, 0) node[left] {$a$};
\end{circuitikz}

\subsubsection{OR Using NOR}

\begin{circuitikz}
  \draw
    % First NOR gate with inputs a and b
    node[nor port] (nor1) at (0, 0) {}
    (nor1.in 1) node[left] {a}
    (nor1.in 2) node[left] {b}

    % Second NOR gate wired as NOT
    node[nor port] (nor2) at (2, 0) {}

    % Tie both inputs of nor2 together
    (nor2.in 1) coordinate (nor2top)
    (nor2.in 2) coordinate (nor2bot)
    (nor2top) -- (nor2bot)

    % Connect nor1 output to the midpoint of nor2's tied inputs
    (nor1.out) -- ($(nor2top)!0.5!(nor2bot)$)

    % Output label
    (nor2.out) node[right] {$a \vee b$};
\end{circuitikz}

\subsubsection{AND Using NOR}

\begin{circuitikz}
  \draw
    % Top NOR gate acting as NOT b
    node[nor port] (nor1) at (0, 1) {}
    (nor1.in 1) coordinate (nor1top)
    (nor1.in 2) coordinate (nor1bot)
    (nor1top) -- (nor1bot)
    ($(nor1top)!0.5!(nor1bot)$) -- ++(-0.5, 0) node[left] {a}

    % Bottom NOR gate acting as NOT a
    node[nor port] (nor2) at (0, -1) {}
    (nor2.in 1) coordinate (nor2top)
    (nor2.in 2) coordinate (nor2bot)
    (nor2top) -- (nor2bot)
    ($(nor2top)!0.5!(nor2bot)$) -- ++(-0.5, 0) node[left] {b}

    % Final NOR gate
    node[nor port] (nor3) at (2, 0) {}

    % Connect NOT gates to final NOR gate
    (nor1.out) -| (nor3.in 1)
    (nor2.out) -| (nor3.in 2)

    % Output
    (nor3.out) node[right] {$a \wedge b$};
\end{circuitikz}


\subsection{NAND Gate}

\subsubsection{NOT Using NAND}

\begin{circuitikz}
  \draw
    node[nand port] (nand) {}

    % Label the output
    (nand.out) node[right] {$\negt a$}

    % Tie both inputs together with a wire and label
    (nand.in 1) coordinate (top)
    (nand.in 2) coordinate (bot)
    (top) -- (bot)

    % Single input wire from the midpoint
    ($(top)!0.5!(bot)$) -- ++(-0.5, 0) node[left] {$a$};
\end{circuitikz}

\subsubsection{OR Using NAND}

\begin{circuitikz}
  \draw
    % Top NAND gate acting as NOT a
    node[nand port] (nand1) at (0, 1) {}
    (nand1.in 1) coordinate (nand1top)
    (nand1.in 2) coordinate (nand1bot)
    (nand1top) -- (nand1bot)
    ($(nand1top)!0.5!(nand1bot)$) -- ++(-0.5, 0) node[left] {a}

    % Bottom NAND gate acting as NOT b
    node[nand port] (nand2) at (0, -1) {}
    (nand2.in 1) coordinate (nand2top)
    (nand2.in 2) coordinate (nand2bot)
    (nand2top) -- (nand2bot)
    ($(nand2top)!0.5!(nand2bot)$) -- ++(-0.5, 0) node[left] {b}

    % Final NAND gate
    node[nand port] (nand3) at (2, 0) {}

    % Connect NOT gates to final NAND gate
    (nand1.out) -| (nand3.in 1)
    (nand2.out) -| (nand3.in 2)

    % Output
    (nand3.out) node[right] {$a \vee b$};
\end{circuitikz}


\subsubsection{AND Using NAND}

\begin{circuitikz}
  \draw
    % First NAND gate with inputs a and b
    node[nand port] (nand1) at (0, 0) {}
    (nand1.in 1) node[left] {a}
    (nand1.in 2) node[left] {b}

    % Second NAND gate wired as NOT
    node[nand port] (nand2) at (2, 0) {}

    % Tie both inputs of nand2 together
    (nand2.in 1) coordinate (nand2top)
    (nand2.in 2) coordinate (nand2bot)
    (nand2top) -- (nand2bot)

    % Connect nand1 output to the midpoint of nand2's tied inputs
    (nand1.out) -- ($(nand2top)!0.5!(nand2bot)$)

    % Output label
    (nand2.out) node[right] {$a \wedge b$};
\end{circuitikz}


\section{Induction and Recursion}
\label{sec:induction-and-recursion}

\subsection{Proof by Mathematical Induction}

To prove that a property $P(n)$ is \textbf{true} for every integer $n \geq b$:
\begin{enumerate}
  \item \textbf{Base Case}: Show that $P(b)$ is true
  \item \textbf{Inductive Hypothesis}: Assume $P(k)$ true for any $k \geq b$
  \item \textbf{Inductive Step}: \textbf{Show} that $P(k) \to P(k+1)$ \textbf{is true}
  \item \textbf{Conclusion}: If $P(b)$ is true and $P(k) \to P(k+1)$ is true, conclude that $P(n)$ is true for all $n \geq b$
\end{enumerate}

\subsection{Proof by Strong Mathematical Induction}
\begin{itemize}
  \item Prove recursive problems with several base cases (e.g., recurrence relations with order larger than 1)
  \item Makes a stronger hypothesis than the principle of Mathematical Induction:
  \begin{itemize}
    \item The \textbf{base step} may contain proofs for \textbf{several initial values}, rather than just the base case
    \item In the inductive step, the trust of $P(k)$ is assumed not just for an arbitrary $k \geq b$ but for \textbf{all values through $k$}
  \end{itemize}
\end{itemize}

Let $b, b$ be two integers with $a \leq b$. To prove that a property $P(n)$ is \textbf{true} for every integer $n \geq a$:
\begin{enumerate}
  \item \textbf{Base Case}: Show that $P(n)$ is true for all base cases from $a$ to $b$, that is: $P(a) \wedge P(a+1) \wedge \dots \wedge P(b)$
  \item \textbf{Inductive Hypothesis}: For any $k geq b$, assume \textbf{$P(i)$ is true for} $i = a, \dots, b, \dots, k$
  \item \textbf{Show} that $P(k+1)$ \textbf{is true}
  \item \textbf{Conclusion}: If 1 and 2 are verified, conclude that $P(n)$ is true for all $n \geq a$
\end{enumerate}

\subsection{Method of the Loop Invariant to Prove Algorithm Correctness}

For every $n \geq 0$, let $P(n)$ be the loop invariant for a loop. \textbf{If the following four properties are true, then the loop is correct} with respect to its pre- and post-conditions:

\begin{enumerate}
  \item \textbf{Basic property}: the loop invariant is true before the loop starts, i.e., $P(0)$ is true before the loop starts. ($P(0)$ is $P(n)$ when $n = 0$). → "The loop invariant is true before the loop starts (i.e., base case is true)"
  \item \textbf{Inductive property}: for any integer $k \geq 0$, if \textbf{both the guard $G$ and the loop invariant $P(k)$} are true before an iteration of the loop, then $P(k+1)$ is true after iteration. → "It remains true after every iteration (i.e., $P(k) \to P(k+1)$)"
  \item \textbf{Eventual Falsity of the Guard}: after a finite number of iterations of the loop, the guard $G$ becomes false. → "The loop eventually stops (i.e., that the loop condition eventually turns false)"
  \item \textbf{Correctness of the post-condition}: if $N$ is the number of iterations after which $G$ is false and $P(n)$ is true, then the post-condition is satisfied. → "At the point when the loop ends, you need to verify that the \textbf{loop invariant} is true"
\end{enumerate}

\section{Functions and Convolution}
\label{sec:functions-and-convolution}

\subsection{Laws of Exponents}
If $b$ and $c$ are any positive real numbers and $u$ and $v$ are any real numbers, the following laws of exponents hold true:

\begin{itemize}
  \item $b^{u}b^{v} = b^{u+v}$
  \item $(b^{u})^{v} = b^{uv}$
  \item $\frac{b^{u}}{b^{v}} = b^{u-v}$
  \item $(bc)^{u} = b^{u}c^{u}$
\end{itemize}

\subsection{Laws of Logarithms}
For any positive real numbers $b$, $c$ and $x$ with $b \neq 1$ and $c \neq 1$:

\begin{itemize}
  \item $\log_{b}(xy) = \log_{b}x + \log_{b}y$
  \item $\log_{b}(\frac{x}{y}) = \log_{b}x - \log_{b}y$
  \item $\log_{b}(x^{a}) = a \log_{b}x$
  \item $\log_{c}x = \frac{\log_{b}x}{\log_{b}c}$
\end{itemize}

\subsection{1D Convolution}
$f(n)$, $g(n)$, $y(n)$: $\mathbb{Z} \to \mathbb{R}$ \quad $\forall n \in \mathbb{Z}$

1D convolution between $f$ and $g$, denoted with $y(n) = (f * g)(n) = \sum_{k = -\infty}^{+\infty} f(k) g(n-k)$

\begin{itemize}
  \item $f$: input function
  \item $g$: kernel or filter
\end{itemize}

To convolve two functions $f$ and $g$: $f * g$
\begin{enumerate}
  \item Flip $g$ horizontally ($-k$)
  \item Slide it over $f$ by adjusting $n$ and compute the sum of the products, starting at the first overlap
  \item Continue until the last overlap
  \item Calculate $y(n)$ for using all $n$-values used to slide $g$ over $f$
\end{enumerate}

\textit{Note}: Since the operation is commutative, sometimes it might be more convenient to compute $g * f$

\subsubsection{Convolution Table Approach}
\begin{enumerate}
  \item Create vector representations of $f$ and $g$
  \item Identify \textit{min index} and \textit{max index} → min and max non-zero values on the $x$-axis for $f$ and $g$
  \item Create a table with $f$ as columns from \textit{min index} to \textit{max index} and $g$ as rows from \textit{min index} to \textit{max index}
  \item Multiply the values
  \item Create $y(n)$ by summing the diagonals
  \item $n = 0$ at the diagonal where $\textit{row index} + \textit{column index} = 0$ (and equivalent for other $n$-values)
\end{enumerate}

\subsection{Properties of Convolution}
\begin{itemize}
  \item \textbf{Commutative Law}: $f * g = g * f$
  \item \textbf{Associative Law}: $(f * g) * h = f * (g * h)$
  \item \textbf{Associative Law with Scalar Multiplication}: $(\lambda f) * g = f * (\lambda g) = \lambda (f * g)$
  \item \textbf{Distributive Law}: $f * (g+h) = (f * g) + (f * h)$
  \item \textbf{Convolution with a Dirac}: $f * \delta = f$
\end{itemize}

\subsection{Why We Flip the Function in the Definition of Convolution}
$(f * g)(n) = (g * f)(n) \Rightarrow \sum_{k = -\infty}^{+\infty} f(k) g(n-k) = \sum_{k = -\infty}^{+\infty} g(k) f(n-k)$

\begin{itemize}
  \item Without flipping, convolution would not be commutative
\end{itemize}


\subsection{2D Convolution}
Let $F$, $G$, $Y$ be \textbf{functions} defined on the infinite set $\mathbb{Z}^{2}$: $F(i,j)$, $G(i,j)$, $Y(i,j)$: $\mathbb{Z}^{2} \to \mathbb{R}$ \quad $\forall(i,j) \in \mathbb{Z}^{2}$

We define the 2D convolution between $f$ and $g$: $Y(i,j) = (F * G)(i,j) = \sum_{u = -\infty}^{+\infty} \sum_{v = -\infty}^{+\infty} F(u,v) G(i-u, j-v)$

\begin{itemize}
  \item $f$: input image
  \item $g$: kernel or filter
\end{itemize}
